\documentclass[12pt,a4paper]{article}
\usepackage[margin=2.5cm]{geometry}
\usepackage{amsmath,amssymb,amsfonts}
\usepackage{physics}
\usepackage{enumitem}
\usepackage{fancyhdr}
\usepackage{lastpage}
\usepackage{graphicx}
\usepackage{multicol}
\usepackage{array}
\usepackage{tabularx}
\usepackage{tikz}

% Page style
\pagestyle{fancy}
\fancyhf{}
\renewcommand{\headrulewidth}{0.4pt}
\renewcommand{\footrulewidth}{0.4pt}
\lhead{PHYS10012}
\rhead{Mock Examination}
\cfoot{Page \thepage\ of \pageref{LastPage}}

% Question numbering
\newcounter{questioncounter}
\newcommand{\question}[1]{%
  \stepcounter{questioncounter}%
  \vspace{0.5cm}%
  \noindent\textbf{Question \thequestioncounter.} \hfill [#1 marks]%
  \vspace{0.2cm}%
  \par%
}

% Sub-part environments
\newlist{parts}{enumerate}{2}
\setlist[parts,1]{label=(\alph*),leftmargin=2em,itemsep=0.3cm}
\setlist[parts,2]{label=(\roman*),leftmargin=2em,itemsep=0.2cm}

% Mark allocation in sub-parts
\renewcommand{\marks}[1]{\hfill [#1]}

% MCQ environment
\newcounter{mcqcounter}
\newcommand{\mcq}[1]{%
  \stepcounter{mcqcounter}%
  \vspace{0.3cm}%
  \noindent\textbf{\themcqcounter.} #1 \hfill [2 marks]%
  \vspace{0.1cm}%
}
\newlist{choices}{enumerate}{1}
\setlist[choices]{label=\Alph*.,leftmargin=2.5em,itemsep=0.1cm}

% Dirac notation shortcuts
\renewcommand{\ket}[1]{\left|#1\right\rangle}
\renewcommand{\bra}[1]{\left\langle #1\right|}
\renewcommand{\braket}[2]{\left\langle #1\middle|#2\right\rangle}
\renewcommand{\ketbra}[2]{\left|#1\right\rangle\!\left\langle #2\right|}
\renewcommand{\expect}[1]{\left\langle #1\right\rangle}

\begin{document}

% Title block
\begin{center}
{\Large\textbf{UNIVERSITY OF BRISTOL}}\\[0.3cm]
{\large Faculty of Science}\\[0.5cm]
{\Large\textbf{MOCK EXAMINATION --- Paper 3}}\\[0.3cm]
{\large\textbf{Core Physics I --- Classical, Quantum and Thermal Physics}}\\[0.2cm]
{\large Module Code: PHYS10012}\\[0.5cm]
{\large Duration: 3 Hours}\\[0.3cm]
{\large Total Marks: 100}\\[0.5cm]
\end{center}

\noindent\rule{\textwidth}{0.4pt}

\vspace{0.3cm}
\noindent\textbf{Instructions:}
\begin{itemize}[itemsep=0.1cm]
\item Answer \textbf{ALL} questions.
\item \textbf{Section A} (Oscillations and Waves): 10 multiple-choice questions, 20 marks total.
\item \textbf{Section B} (Quantum Physics): 10 multiple-choice questions, 20 marks total.
\item \textbf{Section C}: 4 long-answer questions, 60 marks total.
\item An approved calculator is permitted.
\item A formula sheet is provided at the end of this paper.
\item Show all working for Section C. Marks will be awarded for clear, logical solutions.
\end{itemize}

\noindent\rule{\textwidth}{0.4pt}
\newpage

% ============================================================
% SECTION A: OSCILLATIONS AND WAVES MCQ
% ============================================================
\section*{Section A: Oscillations and Waves}
\noindent\textit{Answer all 10 questions. Each question is worth 2 marks.}\\[0.3cm]

\setcounter{mcqcounter}{0}

% QUESTION A1: SHM
\mcq{A mass $m$ on a spring of spring constant $k$ oscillates with period $T$. If the mass is replaced by one of mass $2m$ (with the same spring), what is the new period?}
\begin{choices}
\item $T/\sqrt{2}$
\item $T$
\item $\sqrt{2}\,T$
\item $2T$
\end{choices}

% QUESTION A2: Energy in SHM
\mcq{A particle undergoes simple harmonic motion with amplitude $A$. At what displacement from equilibrium is the kinetic energy equal to three times the potential energy?}
\begin{choices}
\item $A/4$
\item $A/\sqrt{3}$
\item $A/2$
\item $A/\sqrt{2}$
\end{choices}

% QUESTION A3: Damped oscillation
\mcq{An underdamped oscillator has quality factor $Q = 50$. Approximately how many complete oscillation cycles does it take for the amplitude to fall to $1/e$ of its initial value?}
\begin{choices}
\item $\pi/Q \approx 0.063$ cycles
\item $Q/(2\pi) \approx 8$ cycles
\item $Q/\pi \approx 16$ cycles
\item $2\pi Q \approx 314$ cycles
\end{choices}

% QUESTION A4: Forced oscillation
\mcq{A driven damped oscillator has natural frequency $\omega_0$ and damping constant $\gamma$. At what driving frequency $\omega$ is the phase lag $\delta$ between the driving force and the displacement exactly $\pi/2$?}
\begin{choices}
\item $\omega = \omega_0/2$
\item $\omega = \omega_0$
\item $\omega = \omega_0\sqrt{1-\gamma^2/(2\omega_0^2)}$
\item $\omega = \sqrt{\omega_0^2 + \gamma^2}$
\end{choices}

% QUESTION A5: Transverse wave at fixed end
\mcq{A transverse pulse travels along a taut string and arrives at an end that is rigidly fixed to a wall. Which of the following best describes what happens?}
\begin{choices}
\item The pulse is transmitted into the wall with no reflection.
\item The pulse is reflected with no change in phase (upright).
\item The pulse is reflected with a phase change of $\pi$ (inverted).
\item The pulse is absorbed and the string returns to rest.
\end{choices}

% QUESTION A6: String impedance
\mcq{A string has a linear mass density $\mu = 0.04\;\text{kg/m}$ and is held under a tension $T = 100\;\text{N}$. What is the impedance $Z$ of the string?}
\begin{choices}
\item $0.8\;\text{kg/s}$
\item $2.0\;\text{kg/s}$
\item $4.0\;\text{kg/s}$
\item $50\;\text{kg/s}$
\end{choices}

% QUESTION A7: Decibels
\mcq{If the intensity of a sound wave is doubled, by how many decibels does the sound level increase?}
\begin{choices}
\item $1.5\;\text{dB}$
\item $3.0\;\text{dB}$
\item $6.0\;\text{dB}$
\item $10\;\text{dB}$
\end{choices}

% QUESTION A8: Doppler
\mcq{A source emitting sound at frequency $f = 500\;\text{Hz}$ moves toward a stationary observer at $v_s = 20\;\text{m/s}$. At the same time, the observer moves toward the source at $v_o = 10\;\text{m/s}$. Taking the speed of sound as $v = 340\;\text{m/s}$, what is the frequency heard by the observer?}
\begin{choices}
\item $547\;\text{Hz}$
\item $531\;\text{Hz}$
\item $469\;\text{Hz}$
\item $556\;\text{Hz}$
\end{choices}

% QUESTION A9: Standing waves
\mcq{A string of length $L$ is fixed at both ends and vibrates in its 4th harmonic. How many nodes are present on the string (including the endpoints)?}
\begin{choices}
\item $3$
\item $4$
\item $5$
\item $8$
\end{choices}

% QUESTION A10: Coupled oscillators
\mcq{Two identical simple pendulums of length $L$ and mass $m$ are connected by a light spring of spring constant $k_c$. Which of the following correctly describes the two normal modes of oscillation?}
\begin{choices}
\item Both pendulums swing in phase at frequency $\sqrt{g/L}/(2\pi)$; both swing in antiphase at frequency $\sqrt{g/L + 2k_c/m}/(2\pi)$.
\item Both swing in phase at frequency $\sqrt{g/L + k_c/m}/(2\pi)$; both swing in antiphase at frequency $\sqrt{g/L + k_c/m}/(2\pi)$.
\item Both swing in phase at frequency $\sqrt{g/L + 2k_c/m}/(2\pi)$; both swing in antiphase at frequency $\sqrt{g/L}/(2\pi)$.
\item Both swing in phase at frequency $\sqrt{g/L}/(2\pi)$; both swing in antiphase at frequency $\sqrt{g/L + k_c/m}/(2\pi)$.
\end{choices}

\newpage

% ============================================================
% SECTION B: QUANTUM PHYSICS MCQ
% ============================================================
\section*{Section B: Quantum Physics}
\noindent\textit{Answer all 10 questions. Each question is worth 2 marks.}\\[0.3cm]

\setcounter{mcqcounter}{0}

% QUESTION B1: Dirac notation
\mcq{In Dirac notation, what type of mathematical object is $\braket{a}{b}\braket{c}{d}$?}
\begin{choices}
\item A ket
\item A bra
\item An operator
\item A scalar (complex number)
\end{choices}

% QUESTION B2: State validity
\mcq{Which of the following represents a valid (normalised) quantum state?}
\begin{choices}
\item $\ket{\psi} = \ket{1} + \ket{2}$
\item $\ket{\psi} = \frac{1}{2}\ket{1} + \frac{1}{2}\ket{2}$
\item $\ket{\psi} = \frac{1}{\sqrt{2}}\ket{1} + \frac{i}{\sqrt{2}}\ket{2}$
\item $\ket{\psi} = \frac{1}{\sqrt{2}}\ket{1} + \ket{2}$
\end{choices}

% QUESTION B3: Completeness relation
\mcq{If $\{\ket{n}\}$ is a complete orthonormal basis, what does $\sum_n \ketbra{n}{n}$ equal?}
\begin{choices}
\item $0$
\item The Hamiltonian $\hat{H}$
\item The identity operator $\hat{I}$
\item $\sum_n \braket{n}{n}$
\end{choices}

% QUESTION B4: Hermitian conjugate
\mcq{Given two operators $\hat{A}$ and $\hat{B}$, which of the following equals $(\hat{A}\hat{B})^\dagger$?}
\begin{choices}
\item $\hat{A}^\dagger \hat{B}^\dagger$
\item $\hat{B}^\dagger \hat{A}^\dagger$
\item $\hat{B}\hat{A}$
\item $(\hat{A}\hat{B})^{-1}$
\end{choices}

% QUESTION B5: Eigenvalue
\mcq{What are the eigenvalues of the matrix $M = \begin{pmatrix} 3 & 2 \\ 2 & 0 \end{pmatrix}$?}
\begin{choices}
\item $1$ and $2$
\item $4$ and $-1$
\item $3$ and $0$
\item $5$ and $-2$
\end{choices}

% QUESTION B6: Measurement
\mcq{A system is in state $\ket{\psi}$ and observable $A$ is measured, yielding eigenvalue $a_n$ (which is non-degenerate). Immediately after the measurement, the state of the system is:}
\begin{choices}
\item $\ket{\psi}$ (unchanged)
\item $\hat{A}\ket{\psi}$
\item The eigenstate $\ket{a_n}$ corresponding to eigenvalue $a_n$
\item A uniform superposition of all eigenstates of $\hat{A}$
\end{choices}

% QUESTION B7: Time evolution
\mcq{A system is prepared in the energy eigenstate $\ket{E_1}$ at time $t=0$. What is its state at time $t$?}
\begin{choices}
\item $\ket{E_1}$
\item $e^{+iE_1 t/\hbar}\ket{E_1}$
\item $e^{-iE_1 t/\hbar}\ket{E_1}$
\item $\cos(E_1 t/\hbar)\ket{E_1}$
\end{choices}

% QUESTION B8: Spin-1/2
\mcq{For a spin-$\frac{1}{2}$ particle, what are the eigenvalues of the operator $\hat{S}^2 = \hat{S}_x^2 + \hat{S}_y^2 + \hat{S}_z^2$?}
\begin{choices}
\item $\pm\hbar/2$
\item $\hbar^2/4$
\item $\frac{3}{4}\hbar^2$
\item $\hbar^2$
\end{choices}

% QUESTION B9: Fermions
\mcq{Two identical fermions can each occupy one of two single-particle states $\ket{\alpha}$ or $\ket{\beta}$. How many distinct two-particle states are allowed?}
\begin{choices}
\item $1$
\item $2$
\item $3$
\item $4$
\end{choices}

% QUESTION B10: Quarks
\mcq{What is the quark composition of a proton?}
\begin{choices}
\item $udd$
\item $uud$
\item $uds$
\item $u\bar{d}$
\end{choices}

\newpage

% ============================================================
% SECTION C: LONG ANSWER QUESTIONS
% ============================================================
\section*{Section C: Long Answer Questions}
\noindent\textit{Answer all 4 questions. Show all working.}

\setcounter{questioncounter}{0}

% ---- QUESTION C1: MECHANICS (15 marks) ----
\question{15}

\noindent\textbf{Collisions, Impulse, and Centre of Mass}

\begin{parts}
\item A 4\,kg trolley moving at 6\,m/s to the right collides head-on with a 2\,kg trolley moving at 3\,m/s to the left. The collision is perfectly inelastic (the trolleys stick together). Find the velocity of the combined mass after the collision, and state its direction of motion. \qmarks{3}

\item Calculate the total kinetic energy before and after the collision in part~(a). What fraction of the initial kinetic energy is lost? \qmarks{3}

\item Suppose instead that the collision described in part~(a) is \emph{elastic}. Using conservation of momentum and kinetic energy (or the elastic collision formula provided on the formula sheet), find the velocity of each trolley after the collision. \qmarks{4}

\item A constant force acts on a 3\,kg object for 0.5\,s, changing its velocity from 2\,m/s to 8\,m/s (in the same direction). Calculate the impulse delivered to the object and the average force exerted. \qmarks{2}

\item Three point particles of masses 1\,kg, 2\,kg, and 3\,kg are located at positions $(0,\,0)$, $(3,\,0)$, and $(0,\,4)$ metres, respectively. Find the position of the centre of mass of the system. \qmarks{3}
\end{parts}

\newpage

% ---- QUESTION C2: PROPERTIES OF MATTER (15 marks) ----
\question{15}

\noindent\textbf{Thermodynamic Cycle}

\noindent Two moles of monatomic ideal gas ($\gamma = 5/3$) begin in state~A at pressure $P_A = 1.00 \times 10^5\;\text{Pa}$ and temperature $T_A = 300\;\text{K}$.

\begin{parts}
\item Find the volume $V_A$ of the gas in state~A. \qmarks{3}

\item The gas is heated at constant pressure (isobaric process) until the temperature reaches $T_B = 600\;\text{K}$ (state~B). Find:
\begin{parts}
\item the volume $V_B$;
\item the work done by the gas;
\item the heat absorbed by the gas.
\end{parts}
\qmarks{4}

\item From state~B, the gas expands adiabatically until its temperature returns to $300\;\text{K}$ (state~C). Find the volume $V_C$ and the pressure $P_C$. \qmarks{4}

\item The gas is then compressed isothermally at $300\;\text{K}$ back to its original state~A.
\begin{parts}
\item Calculate the work done on the gas in this step.
\item What is the total work done by the gas in one complete cycle?
\item Determine the efficiency of this cycle and compare it with the Carnot efficiency for the same temperature extremes.
\end{parts}
\qmarks{4}
\end{parts}

\newpage

% ---- QUESTION C3: OSCILLATIONS AND WAVES (15 marks) ----
\question{15}

\noindent\textbf{Coupled Oscillators and Pendulums}

\begin{parts}
\item Two identical masses $m = 0.5\;\text{kg}$ are connected in a line by three identical springs, each with spring constant $k = 40\;\text{N/m}$, in the arrangement:

\begin{center}
wall --- spring --- mass\,1 --- spring --- mass\,2 --- spring --- wall
\end{center}

Let $x_1$ and $x_2$ denote the displacements of mass~1 and mass~2 from their respective equilibrium positions. Write down the equations of motion for the two masses. \qmarks{3}

\item Find the two normal mode angular frequencies. Describe the motion of the two masses in each normal mode. \qmarks{4}

\item At $t = 0$, mass~1 is displaced by $5\;\text{cm}$ from equilibrium and mass~2 is at its equilibrium position. Both masses are released from rest. Express the subsequent motion of each mass as a superposition of the two normal modes. \qmarks{3}

\item A simple pendulum has length $L = 1.2\;\text{m}$. Find its period of small oscillations. \qmarks{2}

\item A thin uniform rod of mass $M = 0.8\;\text{kg}$ and length $\ell = 0.6\;\text{m}$ is pivoted at one end and swings as a physical (compound) pendulum. Find:
\begin{parts}
\item the moment of inertia of the rod about the pivot;
\item the distance from the pivot to the centre of mass;
\item the period of small oscillations.
\end{parts}
\qmarks{3}
\end{parts}

\newpage

% ---- QUESTION C4: QUANTUM PHYSICS (15 marks) ----
\question{15}

\noindent\textbf{Operators, Eigenvalues, and Expectation Values}

\noindent Let $\{\ket{1},\,\ket{2}\}$ be an orthonormal basis. Consider the operator
\[
\hat{A} = 2\ketbra{1}{1} - \ketbra{2}{2} + 2\ketbra{1}{2} + 2\ketbra{2}{1}.
\]

\begin{parts}
\item Show that $\hat{A}$ is Hermitian. Write $\hat{A}$ as a $2\times 2$ matrix in the $\{\ket{1},\,\ket{2}\}$ basis. \qmarks{3}

\item Find the eigenvalues and normalised eigenvectors of $\hat{A}$. \qmarks{4}

\item A system is prepared in the state
\[
\ket{\psi} = \frac{3}{5}\ket{1} + \frac{4}{5}\ket{2}.
\]
If observable $A$ is measured, what are the possible outcomes and their probabilities? Calculate the expectation value $\expect{A}$. \qmarks{4}

\item Suppose the measurement in part~(c) yields the \emph{larger} eigenvalue. The system then evolves under the Hamiltonian
\[
\hat{H} = E_0\ketbra{1}{1} + 2E_0\ketbra{2}{2},
\]
where $E_0 > 0$. Find the state of the system at time~$t$, and determine the probability of again measuring the larger eigenvalue of $A$ at that time. \qmarks{4}
\end{parts}

\newpage

% ============================================================
% FORMULA SHEET
% ============================================================
% EQUATION SHEET — PHYS10012 Core Physics I
% This file is \input{} at the end of each exam paper

\newpage
\section*{Formula Sheet}
\noindent\rule{\textwidth}{0.4pt}

\subsection*{Mathematical Formulae}

\begin{align*}
&\text{Binomial expansion:} & (1+x)^n &\approx 1 + nx + \frac{n(n-1)}{2!}x^2 + \cdots \quad (|x| \ll 1) \\[4pt]
&\text{Euler's formula:} & e^{i\theta} &= \cos\theta + i\sin\theta \\[4pt]
&\text{Trig identities:} & \sin A + \sin B &= 2\cos\!\left(\tfrac{A-B}{2}\right)\sin\!\left(\tfrac{A+B}{2}\right) \\
& & \cos A + \cos B &= 2\cos\!\left(\tfrac{A-B}{2}\right)\cos\!\left(\tfrac{A+B}{2}\right)
\end{align*}

\subsection*{Mechanics}

\begin{align*}
&\text{SUVAT:} & v &= u + at, \quad s = ut + \tfrac{1}{2}at^2, \quad v^2 = u^2 + 2as, \quad s = \tfrac{1}{2}(u+v)t \\[4pt]
&\text{Dot product:} & \mathbf{a}\cdot\mathbf{b} &= |\mathbf{a}||\mathbf{b}|\cos\theta = a_x b_x + a_y b_y + a_z b_z \\[4pt]
&\text{Cross product:} & \mathbf{a}\times\mathbf{b} &= |\mathbf{a}||\mathbf{b}|\sin\theta\;\hat{\mathbf{n}} = \begin{vmatrix}\hat{\mathbf{i}}&\hat{\mathbf{j}}&\hat{\mathbf{k}}\\a_x&a_y&a_z\\b_x&b_y&b_z\end{vmatrix} \\[4pt]
&\text{Rocket equation:} & \Delta v &= u \ln\!\left(\frac{m_0}{m_f}\right) \\[4pt]
&\text{Gravitational PE:} & U &= -\frac{GMm}{r} \\[4pt]
&\text{Escape velocity:} & v_e &= \sqrt{\frac{2GM}{R}} \\[4pt]
&\text{Force from PE:} & F &= -\frac{dU}{dx} \\[4pt]
&\text{Elastic collision (1D):} & v_1' &= \frac{m_1 - m_2}{m_1 + m_2}v_1 + \frac{2m_2}{m_1+m_2}v_2 \\[4pt]
&\text{Centre of mass:} & \mathbf{R}_\text{cm} &= \frac{\sum m_i\mathbf{r}_i}{\sum m_i} \\[4pt]
&\text{Reduced mass:} & \mu &= \frac{m_1 m_2}{m_1 + m_2} \\[4pt]
&\text{Centripetal accel:} & a_c &= \frac{v^2}{r} = \omega^2 r \\[4pt]
&\text{Parallel axis thm:} & I &= I_\text{cm} + Md^2 \\[4pt]
&\text{Atwood machine:} & a &= \frac{(m_1-m_2)g}{m_1+m_2}, \quad T = \frac{2m_1 m_2 g}{m_1+m_2}
\end{align*}

\noindent\textbf{Moments of inertia} (about axis through centre of mass):
\begin{center}
\begin{tabular}{ll}
Thin ring (radius $R$): & $I = MR^2$ \\
Solid disk (radius $R$): & $I = \tfrac{1}{2}MR^2$ \\
Solid sphere (radius $R$): & $I = \tfrac{2}{5}MR^2$ \\
Hollow sphere (radius $R$): & $I = \tfrac{2}{3}MR^2$ \\
Thin rod (length $L$, about centre): & $I = \tfrac{1}{12}ML^2$ \\
Thin rod (length $L$, about end): & $I = \tfrac{1}{3}ML^2$ \\
\end{tabular}
\end{center}

\subsection*{Properties of Matter}

\begin{align*}
&\text{Ideal gas:} & PV &= nRT \\[4pt]
&\text{Heat capacities:} & C_P - C_V &= R \quad\text{(per mole)} \\[4pt]
&\text{Adiabatic:} & PV^\gamma &= \text{const}, \quad TV^{\gamma-1} = \text{const} \\[4pt]
&\text{Isothermal work:} & W &= nRT\ln\!\left(\frac{V_2}{V_1}\right) \\[4pt]
&\text{Carnot efficiency:} & \eta &= 1 - \frac{T_c}{T_h} \\[4pt]
&\text{Entropy:} & dS &= \frac{dQ_\text{rev}}{T}, \quad \Delta S = mc\ln\!\left(\frac{T_2}{T_1}\right), \quad \Delta S = nR\ln\!\left(\frac{V_2}{V_1}\right) \\[4pt]
&\text{Lennard-Jones:} & V(r) &= 4\varepsilon\left[\left(\frac{a}{r}\right)^{12} - \left(\frac{a}{r}\right)^{6}\right] \\[4pt]
&\text{Mean KE:} & \langle KE \rangle &= \tfrac{3}{2}k_B T \\[4pt]
&\text{MB speeds:} & c_\text{mp} &= \sqrt{\frac{2k_BT}{m}}, \quad c_\text{avg} = \sqrt{\frac{8k_BT}{\pi m}}, \quad c_\text{rms} = \sqrt{\frac{3k_BT}{m}}
\end{align*}

\subsection*{Oscillations and Waves}

\begin{align*}
&\text{SHM:} & x(t) &= A\cos(\omega_0 t + \phi), \quad \omega_0 = \sqrt{k/m} \\[4pt]
&\text{Damped SHM:} & x(t) &= Ae^{-\gamma t/2}\cos(\omega_d t + \phi), \quad \omega_d = \sqrt{\omega_0^2 - \gamma^2/4} \\[4pt]
&\text{Quality factor:} & Q &= \omega_0/\gamma \\[4pt]
&\text{Forced amplitude:} & A(\omega) &= \frac{F_0/m}{\sqrt{(\omega_0^2 - \omega^2)^2 + \gamma^2\omega^2}} \\[4pt]
&\text{Phase lag:} & \tan\delta &= \frac{\gamma\omega}{\omega_0^2 - \omega^2} \\[4pt]
&\text{Pendulums:} & T_\text{simple} &= 2\pi\sqrt{L/g}, \quad T_\text{physical} = 2\pi\sqrt{I/(Mgh)} \\[4pt]
&\text{Wave equation:} & \frac{\partial^2 y}{\partial x^2} &= \frac{1}{v^2}\frac{\partial^2 y}{\partial t^2} \\[4pt]
&\text{String speed:} & v &= \sqrt{T/\mu} \\[4pt]
&\text{Sound in gas:} & v &= \sqrt{\gamma P/\rho} \\[4pt]
&\text{Harmonic wave:} & y &= A\sin(kx - \omega t) \\[4pt]
&\text{String impedance:} & Z &= \mu v = \sqrt{\mu T} \\[4pt]
&\text{Reflection:} & r &= \frac{Z_1 - Z_2}{Z_1 + Z_2}, \quad t = \frac{2Z_1}{Z_1+Z_2} \\[4pt]
&\text{Power (string):} & \langle P\rangle &= \tfrac{1}{2}\mu\omega^2 A^2 v \\[4pt]
&\text{Decibels:} & \beta &= 10\log_{10}(I/I_0), \quad I_0 = 10^{-12}\;\text{W/m}^2 \\[4pt]
&\text{Doppler:} & f' &= f\frac{v \pm v_o}{v \mp v_s} \\[4pt]
&\text{Group velocity:} & v_g &= \frac{d\omega}{dk} \\[4pt]
&\text{String harmonics:} & f_n &= \frac{nv}{2L} \quad (n=1,2,3,\ldots) \\[4pt]
&\text{Open-closed pipe:} & f_n &= \frac{nv}{4L} \quad (n=1,3,5,\ldots)
\end{align*}

\subsection*{Quantum Physics}

\begin{align*}
&\text{Schr\"{o}dinger eq:} & i\hbar\frac{d}{dt}\ket{\psi(t)} &= H\ket{\psi(t)} \\[4pt]
&\text{Time evolution:} & \ket{\psi(t)} &= \sum_n c_n e^{-iE_n t/\hbar}\ket{E_n}, \quad U(t) = e^{-iHt/\hbar} \\[4pt]
&\text{Spin operators:} & S_z &= \frac{\hbar}{2}\bigl(\ketbra{\!\uparrow}{\uparrow\!} - \ketbra{\!\downarrow}{\downarrow\!}\bigr) \\
& & S_x &= \frac{\hbar}{2}\bigl(\ketbra{\!\uparrow}{\downarrow\!} + \ketbra{\!\downarrow}{\uparrow\!}\bigr) \\
& & S_y &= \frac{\hbar}{2}\bigl(-i\ketbra{\!\uparrow}{\downarrow\!} + i\ketbra{\!\downarrow}{\uparrow\!}\bigr) \\[4pt]
&\text{Spin eigenstates:} & \ket{\!\uparrow_x} &= \tfrac{1}{\sqrt{2}}\bigl(\ket{\!\uparrow_z} + \ket{\!\downarrow_z}\bigr), \quad \ket{\!\downarrow_x} = \tfrac{1}{\sqrt{2}}\bigl(\ket{\!\uparrow_z} - \ket{\!\downarrow_z}\bigr) \\
& & \ket{\!\uparrow_y} &= \tfrac{1}{\sqrt{2}}\bigl(\ket{\!\uparrow_z} + i\ket{\!\downarrow_z}\bigr), \quad \ket{\!\downarrow_y} = \tfrac{1}{\sqrt{2}}\bigl(\ket{\!\uparrow_z} - i\ket{\!\downarrow_z}\bigr) \\[4pt]
&\text{Pauli matrices:} & \sigma_x &= \begin{pmatrix}0&1\\1&0\end{pmatrix}, \quad \sigma_y = \begin{pmatrix}0&-i\\i&0\end{pmatrix}, \quad \sigma_z = \begin{pmatrix}1&0\\0&-1\end{pmatrix}
\end{align*}

\subsection*{Physical Constants}

\begin{center}
\begin{tabular}{lll}
Speed of light & $c$ & $3.00 \times 10^{8}$ m/s \\
Planck's constant & $h$ & $6.626 \times 10^{-34}$ J$\cdot$s \\
Reduced Planck's constant & $\hbar$ & $1.055 \times 10^{-34}$ J$\cdot$s \\
Boltzmann constant & $k_B$ & $1.381 \times 10^{-23}$ J/K \\
Gas constant & $R$ & $8.314$ J/(mol$\cdot$K) \\
Avogadro's number & $N_A$ & $6.022 \times 10^{23}$ mol$^{-1}$ \\
Gravitational constant & $G$ & $6.674 \times 10^{-11}$ N$\cdot$m$^2$/kg$^2$ \\
Acceleration due to gravity & $g$ & $9.81$ m/s$^2$ \\
Electron mass & $m_e$ & $9.109 \times 10^{-31}$ kg \\
Proton mass & $m_p$ & $1.673 \times 10^{-27}$ kg \\
Elementary charge & $e$ & $1.602 \times 10^{-19}$ C \\
Mass of Earth & $M_\oplus$ & $5.97 \times 10^{24}$ kg \\
Radius of Earth & $R_\oplus$ & $6.37 \times 10^{6}$ m \\
\end{tabular}
\end{center}

\noindent\rule{\textwidth}{0.4pt}


\end{document}
